\chapter{Requisitos y arquitectura de la solución}\label{cap:arquitectura}

\section{Requisitos funcionales}\label{sec:requisitos-funcionales}

Los requisitos funcionales identifican capacidades observables que la plataforma debe ofrecer al usuario operador y al evaluador externo. La tabla~\ref{tab:requisitos-funcionales} los enumera con identificador, descripción y criterio de aceptación verificable.

\begin{table}[h]
\footnotesize
\centering
\renewcommand{\arraystretch}{1.15}
\begin{tabular}{|L{1.1cm}|L{5.2cm}|L{6.2cm}|}
\hline
\rowcolor{tema!10}
\bft{ID} & \bft{Requisito} & \bft{Criterio de aceptación} \\ \hline
RF-01 & Despliegue reproducible de OpenMetadata y sus dependencias. & El script \texttt{launch\_infra.ps1} levanta el clúster con \texttt{kubectl get pods} sin errores. \\ \hline
RF-02 & PostgreSQL de referencia con capas \texttt{bronze}/\texttt{silver}/\texttt{gold}. & \texttt{opendata\_demo\_init.sql} aplicado; los tres esquemas son visibles. \\ \hline
RF-03 & Doble ingesta de PostgreSQL en OpenMetadata. & Dos \texttt{DatabaseService} diferenciados con tablas \texttt{gold}. \\ \hline
RF-04 & Selección de datasets publicables. & Solo filas con \texttt{publicar=si} aparecen en el catálogo exportado. \\ \hline
RF-05 & Curación funcional por \ac{FQN} de tabla. & La hoja \ac{CSV} identifica filas únicas por \texttt{table\_fqn}. \\ \hline
RF-06 & Sincronización idempotente con OpenMetadata. & Una segunda ejecución consecutiva aplica 0 cambios. \\ \hline
RF-07 & Exportación a \ac{JSONLD} conforme al perfil activo. & Se genera \texttt{dcat\_catalog.jsonld} con clases \ac{DCAT} activas. \\ \hline
RF-08 & Validación \ac{SHACL} contra el bundle vendorizado. & Se genera \texttt{dcat\_validation\_report.ttl} sin violaciones bloqueantes. \\ \hline
RF-09 & Consola web operativa para todo el flujo. & La \ac{UI} ejecuta cada etapa y muestra logs y artefactos. \\ \hline
\end{tabular}
\caption{Requisitos funcionales con criterio de aceptación. Elaboración propia.}
\label{tab:requisitos-funcionales}
\end{table}

\section{Requisitos no funcionales}\label{sec:requisitos-no-funcionales}

Los requisitos no funcionales recogen las propiedades de calidad transversales que la plataforma debe sostener (tabla~\ref{tab:requisitos-no-funcionales}). Se priorizan reproducibilidad e idempotencia por encima de capacidades habituales en sistemas productivos (alta disponibilidad, \ac{RBAC} avanzado, escalado), que se declaran fuera de alcance en la sección~\ref{sec:no-objetivos}.

\begin{table}[h]
\footnotesize
\centering
\renewcommand{\arraystretch}{1.15}
\begin{tabular}{|L{2.8cm}|L{9.7cm}|}
\hline
\rowcolor{tema!10}
\bft{Atributo} & \bft{Compromiso} \\ \hline
Reproducibilidad & Cualquier ejecución parte de comandos versionados; el repositorio limpio reconstruye el entorno completo. \\ \hline
Idempotencia & La sincronización contra \ac{OM} no duplica etiquetas, propiedades ni recursos. \\ \hline
Simplicidad operativa & Una única capa Python concentra reglas de negocio; \ac{CLI}, scripts y web son envoltorios. \\ \hline
Trazabilidad & Cada artefacto se asocia a configuración, código y comando de origen. \\ \hline
Separación de responsabilidades & La consola web no implementa reglas; las delega en el núcleo. \\ \hline
Portabilidad & El despliegue es transferible entre entornos compatibles con Kubernetes y Helm. \\ \hline
Auditabilidad & Cada ejecución produce evidencia regenerable y consultable. \\ \hline
\end{tabular}
\caption{Requisitos no funcionales y compromisos asociados. Elaboración propia.}
\label{tab:requisitos-no-funcionales}
\end{table}

\section{No objetivos}\label{sec:no-objetivos}

Quedan deliberadamente fuera del alcance funcional la alta disponibilidad, el hardening avanzado, la integración con \ac{SSO}/\ac{LDAP}, el \ac{RBAC} granular, los backups productivos, la observabilidad avanzada, el escalado horizontal y la cosecha activa desde CKAN. Estas exclusiones no son carencias de implementación: corresponden a decisiones de diseño que evitan que una capa adicional desplace la atención del problema principal, que es el gobierno reproducible y la conformidad \ac{DCATAPES}.

\section{Arquitectura lógica}

La arquitectura se organiza en cinco capas. En lugar de describirlas con una tabla, se representan mediante un diagrama (figura~\ref{fig:arq-logica-detalle}), que hace explícita la responsabilidad de cada capa y la dirección de las dependencias y del flujo de datos entre ellas~\cite{mermaidDocs}. De abajo arriba, la fuente técnica (PostgreSQL de referencia) alimenta el catálogo técnico (OpenMetadata), del que el gobierno funcional toma los activos descubiertos; el núcleo Python fusiona hoja y defaults, aplica gobierno idempotente sobre OpenMetadata y produce el catálogo \ac{DCATAPES} validado; la capa de operación (\ac{CLI}, scripts y consola web) invoca ese mismo núcleo sin duplicar reglas.

Cada capa se respalda con evidencia versionada en el repositorio: la fuente técnica en \fqn{sql/opendata_demo_init.sql}; el catálogo técnico en \fqn{ingest_postgres_double.ps1}; el gobierno funcional en \fqn{gold_governance.csv} y \fqn{governance_defaults.yaml}; el núcleo en \fqn{tfm_ingestor/src/tfm_ingestor/}; y la operación en \fqn{scripts/infra/} y \fqn{web/}.

\figuraMemoria{fig_arquitectura_logica_detalle.png}{Arquitectura lógica por capas y responsabilidades: fuente técnica, catálogo técnico, gobierno funcional, núcleo Python y operación. Elaboración propia.}{fig:arq-logica-detalle}

\section{Arquitectura física}

La vista física (figura~\ref{fig:arq-fisica}) muestra cómo se materializan esas capas en el despliegue. Todo se ejecuta sobre un único host —portátil, VPS o nube— con Docker como motor de contenedores y un clúster Kubernetes local gestionado por Kind. OpenMetadata, sus dependencias (MySQL y OpenSearch) y el PostgreSQL de referencia se ejecutan como cargas dentro del clúster. El operario interactúa a través de la consola web y del \ac{CLI}, que se comunican con OpenMetadata mediante un \texttt{port-forward} sobre la \ac{API} \ac{REST} autenticada con \ac{JWT}. El detalle de los \emph{values} Helm y de los manifiestos se recoge en el Anexo~\ref{ap:infraestructura}.

\figuraMemoria{fig_arquitectura_fisica.png}{Arquitectura física de despliegue: host con Docker y clúster Kind que ejecuta OpenMetadata, sus dependencias y el PostgreSQL de referencia. Elaboración propia.}{fig:arq-fisica}

\section{Flujo extremo a extremo}

El flujo implementado sigue una secuencia cerrada. Primero se levanta la infraestructura. Después se ingiere PostgreSQL dos veces en OpenMetadata, bajo dos servicios diferenciados. A continuación se preparan etiquetas y custom properties. El workflow refresca la hoja de gobierno, permite revisar los metadatos funcionales, aplica cambios en OpenMetadata, exporta JSON-LD y ejecuta la validación SHACL.

\figuraMemoria{fig_flujo_operativo_validacion.png}{Flujo operativo del caso de uso de validación. Elaboración propia.}{fig:flujo-operativo}

La figura \ref{fig:pipeline-metadatos} muestra la transformación conceptual desde activos técnicos hasta salida validada.

\figuraMemoria{fig_pipeline_metadatos.png}{Pipeline de metadatos desde OpenMetadata hasta validación SHACL. Elaboración propia.}{fig:pipeline-metadatos}

\section{Decisiones de diseño}

\subsection{PostgreSQL como fuente canónica}

El origen técnico se fija en PostgreSQL porque permite controlar la estructura, reproducir el entorno y demostrar descubrimiento real de activos técnicos. El SQL de referencia crea seis tablas distribuidas en tres capas. Las tablas \fqn{gold.movilidad_resumen_municipio} y \fqn{gold.agenda_cultural_publica} representan el catálogo publicable.

\subsection{Gobierno restringido a la capa \texttt{gold}}\label{sec:decision-gold-only}

La arquitectura \emph{medallion} (bronze/silver/gold) es una convención consolidada en plataformas de datos modernas, popularizada por Databricks como patrón de zonas progresivas de refinamiento~\cite{databricksMedallion}, que separa tres estados de madurez: datos crudos ingeridos sin transformar (\texttt{bronze}), datos limpios y conformados (\texttt{silver}) y datos refinados, agregados y listos para consumo final (\texttt{gold}). La plataforma gobierna \emph{únicamente} la capa \texttt{gold}; las capas \texttt{bronze} y \texttt{silver} se ingieren en \ac{OM} pero no se exportan como datasets \ac{DCATAPES}. Esta restricción es deliberada y se apoya en cuatro razones convergentes.

\paragraph{Conformidad semántica con el perfil DCAT-AP-ES.} El perfil exige que un \texttt{dcat:Dataset} represente un conjunto de datos «publicable», con título, descripción, publicador, temática y, en el caso \ac{HVD}, categoría y legislación aplicables. Las tablas \texttt{bronze} y \texttt{silver} no satisfacen ese contrato: contienen estructuras transitorias, esquemas en evolución o columnas técnicas de ingesta que no tienen significado funcional fuera del pipeline interno. Publicarlas como datasets falsearía el catálogo y haría fracasar la validación \ac{SHACL} en propiedades obligatorias.

\paragraph{Coherencia con el marco DAMA-DMBOK.} \ac{DMBOK} establece que el gobierno aplica sobre activos cuyo significado, propiedad y calidad están declarados y son estables~\cite{damaDmbok}. Las capas \texttt{bronze} y \texttt{silver} son explícitamente transitorias en la arquitectura medallion y no cumplen ese requisito; gobernarlas obligaría a establecer responsables, descripciones y temáticas para datos cuya forma no es estable, generando deuda editorial sin contrapartida.

\paragraph{Separación entre catálogo técnico y catálogo de publicación.} \ac{OM} sigue actuando como inventario técnico completo: las tres capas son visibles, navegables y trazables desde su \ac{UI}, lo que mantiene la utilidad de \ac{OM} como herramienta interna. La frontera se aplica únicamente en la exportación a \ac{DCATAPES}: el catálogo de publicación es un subconjunto curado del catálogo técnico, no su réplica. Esta separación reproduce el comportamiento real de los portales de datos abiertos como datos.gob.es o data.europa.eu, que nunca publican el inventario técnico completo de una organización.

\paragraph{Foco en el valor evaluado por el tribunal.} El alcance funcional declarado en la sección~\ref{sec:requisitos-funcionales} cubre el ciclo completo de gobierno hasta validación \ac{SHACL}. Añadir \texttt{bronze} y \texttt{silver} al gobierno multiplicaría el trabajo editorial sin añadir capacidad demostrable: el flujo, los mecanismos de idempotencia y las shapes \ac{SHACL} son los mismos. La restricción concentra el esfuerzo en demostrar el ciclo completo sobre un subconjunto reducido y representativo, no en escalarlo sobre datos sin valor editorial.

La decisión, por tanto, no es una limitación de implementación sino una alineación con la semántica del perfil destino, con el marco de gobierno aplicado y con la práctica habitual de los portales de datos abiertos. La plataforma puede gobernar tablas adicionales sin cambios estructurales: basta con marcar \fqn{publicar=si} en la hoja \fqn{gold_governance.csv} para los datasets nuevos, lo que demuestra que la decisión es de alcance, no de capacidad técnica.

\subsection{Doble ingesta como validación multi-fuente}

La doble ingesta usa la misma base de referencia, pero crea dos \texttt{DatabaseService} en OpenMetadata: \texttt{postgres\_demo\_service} y \texttt{postgres\_validation\_service}. El objetivo no es duplicar datos por volumen, sino comprobar que el gobierno funciona cuando existen tablas con el mismo nombre en servicios distintos. Por eso la hoja funcional conserva \texttt{schema\_name}, \texttt{table\_name} y, sobre todo, \texttt{table\_fqn}.

\subsection{Hoja funcional para gobierno editorial}

En este trabajo, \emph{curación} designa la actividad editorial de seleccionar, describir, enriquecer y mantener los metadatos de un activo para que sean comprensibles y publicables; es el término consolidado en gestión del dato (\emph{data curation}~\cite{damaDmbok}), frente a «edición» —demasiado genérica— o a «limpieza» y «saneamiento», que aluden a la calidad del dato de negocio y no a sus metadatos. La \emph{curación funcional} es esa misma actividad aplicada a los metadatos funcionales o de negocio de cada dataset (título, descripción, publicador, temática y categoría \ac{HVD}), por oposición a los metadatos técnicos que \ac{OM} captura de forma automática.

La curación funcional se concentra en \texttt{tfm\_ingestor/config/gold\_governance.csv}. Esta hoja permite editar los metadatos que varían por dataset: decisión de publicación, título, descripción, publicador, temática DCAT, categoría HVD y URL de acceso. Los campos técnicos (esquema, tabla y \ac{FQN}) no deben editarse manualmente salvo regeneración controlada del caso de uso. La tabla~\ref{tab:hoja-gobierno-columnas} resume las columnas; el detalle completo con criterios de validación se ofrece en el anexo~\ref{ap:ingesta-gobierno}.

\begin{table}[h]
\footnotesize
\centering
\renewcommand{\arraystretch}{1.15}
\begin{tabular}{|L{4.4cm}|L{1.9cm}|L{5.6cm}|}
\hline
\rowcolor{tema!10}
\bft{Columna} & \bft{Tipo} & \bft{Propósito} \\ \hline
\fqn{publicar}                  & \texttt{si}/\texttt{no} & Inclusión del dataset en el catálogo exportado. \\ \hline
\fqn{schema_name}               & texto       & Esquema PostgreSQL (\texttt{gold} en el caso de uso). \\ \hline
\fqn{table_name}                & texto       & Nombre técnico de la tabla. \\ \hline
\fqn{table_fqn}                 & texto       & Identificador completo en \ac{OM}; clave única. \\ \hline
\fqn{titulo_dataset}            & texto       & Título funcional editorial. \\ \hline
\fqn{descripcion_dataset}       & texto       & Descripción editorial del dataset. \\ \hline
\fqn{publicador}                & texto       & Nombre del organismo publicador. \\ \hline
\fqn{tematica_dcat}             & enum NTI-RISP & Temática \ac{NTIRISP} del dataset. \\ \hline
\fqn{categoria_hvd}             & enum HVD    & Categoría del Reglamento (UE) 2023/138. \\ \hline
\fqn{access_url_distribucion}   & URL         & URL pública de la distribución. \\ \hline
\end{tabular}
\caption{Columnas de la hoja funcional \texttt{gold\_governance.csv}. Elaboración propia.}
\label{tab:hoja-gobierno-columnas}
\end{table}

El campo \fqn{publicar} merece una explicación de gobierno y no de mero filtrado técnico. La exportación del catálogo es idempotente y, por defecto, recorre el total de activos gobernados —datasets, distribuciones y servicios de datos—. El valor \texttt{si}/\texttt{no} introduce un control editorial sobre qué entra realmente en el catálogo \ac{DCATAPES} publicado. Esa distinción alinea la plataforma con el principio \ac{DAMA} de gestionar el ciclo de vida del dato~\cite{damaDmbok}: permite diferenciar los datasets publicables de aquellos que permanecen privados, en estado de borrador o catalogados internamente pero todavía no federables. De este modo, una misma hoja puede reflejar estados de publicación distintos por dataset y acompañar su evolución hacia la federación sin perder ni trazabilidad ni reproducibilidad.

\subsection{Configuración global por YAML}

Los metadatos globales del catálogo se derivan de \texttt{governance\_defaults.yaml}: nombre del catálogo, descripción, publicador, URI de organismo, página principal, idioma, licencia, legislación HVD, contacto y bases de URLs. Estos valores operan como \emph{defaults} a dos niveles: describen el \texttt{dcat:Catalog} y, además, sirven de valor por defecto heredable por cada \texttt{dcat:Dataset}. Esta separación evita repetir información común en cada dataset y reduce errores de consistencia; cuando un dataset concreto requiere un publicador, una licencia, un idioma o un contacto propios —habitual cuando coexisten varios organismos o responsables—, la hoja funcional puede sobrescribir el default global a nivel de fila sin modificar el núcleo. El listado~\ref{code:operational-profile-yaml} muestra el perfil operativo que conecta esos defaults con la hoja \ac{CSV} y fija el caso \ac{HVD} como modo de validación.

\begin{lstlisting}[language=yaml, caption={Perfil operativo \texttt{operational\_profile.yaml} que enlaza defaults, reglas, hoja y caso de validación.}, label=code:operational-profile-yaml, captionpos=b]
defaults_path: "tfm_ingestor/config/governance_defaults.yaml"
rules_path:    "tfm_ingestor/config/mapping_rules.yaml"
sheet_path:    "tfm_ingestor/config/gold_governance.csv"

workflow:
  profile_case:   "hvd"
  allow_warnings: false
  refresh_sheet:  true
  export_output:  "dcat_catalog.jsonld"
  report_output:  "tmp_pytest/dcat_validation_report.ttl"
\end{lstlisting}

\subsection{Núcleo único de lógica}

La regla arquitectónica principal es que la lógica de gobierno vive en Python. El \ac{CLI}, los scripts y la web invocan ese núcleo. Esta decisión evita que la interfaz web acabe definiendo reglas paralelas y facilita probar la solución con \texttt{pytest}. El listado~\ref{code:workflow-service} muestra la firma del entrypoint canónico, una \texttt{dataclass} inmutable (\texttt{frozen=True}) que fija de manera explícita todos los parámetros del workflow: rutas de configuración, conexión a \ac{OM}, modo \emph{dry-run}, refresco de hoja, salidas de exportación y caso de validación. El término \emph{dry-run} (literalmente «ejecución en seco») designa una ejecución que calcula y muestra el plan de cambios sin aplicarlos sobre \ac{OM}. Se conserva en inglés porque es el nombre literal de la opción del \ac{CLI} (\texttt{-{}-dry-run}) y un estándar de facto en herramientas de línea de comandos e \emph{Infrastructure as Code}; traducirlo desincronizaría la memoria del comando real que ejecuta el operario.

\begin{lstlisting}[language=Python, caption={Configuración inmutable del workflow canónico en \texttt{workflow\_service.py}.}, label=code:workflow-service, captionpos=b]
from dataclasses import dataclass

@dataclass(frozen=True)
class WorkflowRunConfig:
    defaults_path: str
    rules_path:    str
    sheet_path:    str
    base_url:      str
    token:         str | None
    limit:         int  = 1000
    dry_run:       bool = False
    refresh_sheet: bool = True
    export_output: str  = "dcat_catalog.jsonld"
    report_output: str | None = None
    allow_warnings: bool = False
    profile_case:  str  = "hvd"
    plan_output:   str | None = None
    tables_input_path: str | None = None
    skip_export:   bool = False
    skip_validate: bool = False


def run_workflow(config: WorkflowRunConfig) -> dict[str, Any]:
    defaults = load_defaults(config.defaults_path)
    rules    = load_rules(config.rules_path)
    api      = OpenMetadataApi(base_url=config.base_url, jwt_token=config.token)
    tables   = discover_tables(api=api, limit=config.limit)
    # 1) refresca hoja  2) sincroniza gobierno  3) exporta DCAT  4) valida SHACL
    ...
\end{lstlisting}

\section{Modelo activo DCAT-AP-ES}

El modelo activo del TFM cubre las clases \texttt{dcat:Catalog}, \texttt{dcat:Dataset}, \texttt{dcat:Distribution}, \texttt{dcat:DataService} y \texttt{foaf:Agent}. OpenMetadata mantiene únicamente los metadatos que conviene gobernar manualmente: \texttt{displayName}, \texttt{description}, \fqn{dcat_publisher_name}, \fqn{dcat_hvd_category}, \fqn{dcat_access_url} y tags \fqn{dcat_theme.*}. El resto se deriva durante la exportación.

\figuraMemoria{fig_mapeo_activo_dcat_openmetadata.png}{Mapeo activo entre clases DCAT-AP-ES y elementos de la plataforma. Elaboración propia.}{fig:mapeo-activo}

Esta decisión prioriza el camino completo sobre la cobertura exhaustiva. El alcance funcional declarado en el capítulo~\ref{cap:objetivos} no incluye los metadatos opcionales que no inciden en la conformidad \ac{HVD}; modelarlos exigiría duplicar trabajo editorial sin aportar valor adicional al ciclo de validación, y queda explícitamente reservado como evolución posterior. La plataforma demuestra, por tanto, un camino completo y validable para el subconjunto obligatorio y operativo del caso \ac{HVD}.
